\section{The shortcut is also an attack surface}
\label{sec:attack}
The source shortcut does not only make detectors over-block benign skills from unseen sources; it lets
an attacker evade them. Appending benign source-style tokens to a malicious skill is a good-word
attack~\cite{lowd2005}; the honest question is whether the \emph{source-style} tokens help beyond the
dilution that appending any text causes to a TF-IDF vector. We answer with a dose-response on the linear
word-LR detector, whose benign-weighted features are directly readable: at each size $k$ we append its
top-$k$ benign-weighted features, and as a control $k$ random vocabulary tokens, to every malicious test
skill it catches at its deployment threshold. The same construction would need a gradient or query budget
against the transformers; we demonstrate the mechanism where the shortcut is legible.

At every $k$ the source-style tokens evade far more than random ones. With only \EvasionK{} injected
tokens, \EvasionAsrAll{} of caught malicious skills fall below threshold, against \EvasionControlAll{}
for the same number of random tokens --- a source-specific gap of \EvasionGapAll{} that padding does
not explain. The attack changes no payload and queries no model.

The mixed-only model of \S\ref{sec:fix} is more robust here too: the same \EvasionK{}-token attack
succeeds on only \EvasionAsrMixed{} of the skills it catches. Because it no longer leans on source
style, the training-set-composition fix that lowers the benign false-positive rate also shrinks the
evasion surface. One bias thus produces both failures --- over-blocking benign skills and letting
malicious ones evade --- and one data fix moves both.

We state two honesty notes. The success rates are over each model's \emph{caught} malicious skills, and
the two models catch different amounts before any attack (at its threshold the all-data model flags
\EvasionBaseRecallAll{} of malicious test skills, the mixed-only model only \EvasionBaseRecallMixed{}), so
the mixed-only model's lower attack rate partly reflects its lower base recall; its source-specific gap
is \EvasionAsrMixed{} benign vs \EvasionControlMixed{} random. The effect sizes are modest: this is a
mechanism demonstration, not a strong attack. It concerns the cheap static pre-install filter, and is
orthogonal to evasions that hide the payload from an LLM judge~\cite{pretext2026,cloakdetonate2026}.
